\section{Introduction}
RGB-D perception is a standard building block for object localisation and robotic
manipulation. Because real robot time is slow, fragile, and expensive, a large share of
this work is developed in simulation and only later moved onto hardware. The recurring
difficulty is that pipelines which behave well in simulation often lose accuracy or fail
outright once they meet a real sensor and a real arm. This discrepancy is known as the
reality gap or sim-to-real gap, and it was already identified in evolutionary robotics
three decades ago, where the role of noise in bridging the gap was made explicit
\cite{jakobi1995noise}. A common modern response is domain randomisation, which
deliberately varies simulator parameters so that the real world looks like one more
sample of the training distribution \cite{tobin2017domain}. Randomisation is effective
but blunt: it rarely tells us which mismatch actually mattered in a given setup.

This thesis isolates a single, measurable source of mismatch in a deliberately
constrained setting. We use a fixed top-down camera over a tabletop, so the dominant
uncertainty is depth sensing rather than viewpoint or hand-eye drift. The camera is a
Time-of-Flight (ToF) device, whose depth error depends on range, surface properties, and
scene geometry. The question we ask is narrow and quantitative: if we measure the real
ToF noise and inject a faithful model of it into simulation, does the simulated
pose-estimation and grasping behaviour move closer to the real system?

The answer we report is, for this pipeline, no, and we argue that this negative result is
the most useful outcome of the project. The perception component is a classical,
outlier-robust estimator, and it turns out to absorb realistic ToF noise almost
completely. Rather than treat that as a failure, we quantify exactly how much noise the
estimator tolerates before it breaks, show that the breaking point sits an order of
magnitude above reality, and confirm on the real robot that the simulated conclusion
holds end to end. The practical message is that, for an estimator of this kind, effort
spent on high-fidelity depth-noise models is misdirected; the residual gap lives
elsewhere, in calibration and lighting.

\subsection{Research Questions}
The work is organised around one main research question and three sub-questions.

\noindent\textbf{Main RQ.} How does modelling realistic Time-of-Flight (ToF) sensor
noise affect the sim-to-real gap in RGB-D tabletop pose estimation and manipulation?
\begin{itemize}
  \item \textbf{SQ1 (noise statistics).} What is the real ToF sensor's noise behaviour,
  and how well do literature models (Nguyen 2012, Kinect-class) describe it? Addressed in
  Phase~2.
  \item \textbf{SQ2 (pose accuracy).} Does injecting measured noise into simulation
  change pose-estimation accuracy relative to clean simulation? Addressed in Phase~3.
  \item \textbf{SQ3 (manipulation propagation).} To what extent do pose estimation errors
  propagate to manipulation performance (pick success rate and placement accuracy)?
  Addressed in Phases~4A (simulation), 4B (real perception), and 4C (real manipulation).
\end{itemize}

\subsection{Hypothesis}
The proposal stated the directional hypothesis that injecting measured ToF noise into
simulation will reduce the sim-to-real gap. As the results show, this hypothesis is not
confirmed at realistic noise levels. We treat the refutation as a finding in its own
right and design the later phases to test whether the same null carries from pose
estimation through to manipulation.

\subsection{Contributions}
\begin{enumerate}
  \item A measured ToF noise characterisation of the OAK-D SR PoE in a fixed top-down
  tabletop setup, including the finding that bulk (body) depth is sub-millimetre and
  essentially flat with range, while edge flying-pixels dominate the pose-relevant error.
  \item A practical depth-noise injection module that implements four literature ToF
  terms (axial Gaussian, lateral jitter, edge mixed-pixels, and dropout) and exposes an
  axial-scale control for ablation.
  \item The central, controlled result: at the measured sensor noise level, the injected
  noise has no measurable effect on pose RMSE or on grasp success for a depth-segmented
  median-centroid estimator. Degradation appears only at roughly ten times the measured
  axial noise for pose and roughly twenty times for manipulation. Because the injector
  demonstrably breaks perception when scaled far enough, the null is sound rather than an
  artefact.
  \item End-to-end validation on the real system: a real-versus-sim planar pose gap
  between $-1.4$ and $+3.6$\,mm across four objects, and a 60-trial real closed-loop
  pick-and-place (34/60 lifts) showing that, once visual feedback corrects the aiming
  error, real grasp success is governed by object grip mechanics rather than by depth
  noise or the simulated pose-error propagation.
  \item A reproducible Gazebo Classic and ROS~2 tabletop manipulation setup for the
  xArm~1S with an overhead RGB-D sensor.
\end{enumerate}
It is important to be precise about scope. The pick-and-place component is the proposal's
authorised predefined grasp strategy for known object types; it is a baseline used to
test whether the pose-level null propagates to grasping, and it is not presented as a
novel system. The engineering of the real closed-loop executor (Phase~4C) is reported as
methods and appendix material, not as a results contribution.

\subsection{Thesis Outline}
Section~\ref{sec:related} reviews related work. Section~\ref{sec:methods} describes the
simulation, the noise characterisation protocol, the injection model, the perception and
grasp pipelines, and the real-world calibration and control. Section~\ref{sec:setup}
specifies the hardware, objects, conditions, trial counts, and metrics.
Section~\ref{sec:results} reports the results. Section~\ref{sec:discussion} interprets the
negative result and its limitations, and Section~\ref{sec:conclusion} concludes and
outlines future work.

% =====================================================================
